\documentclass{article}
\usepackage{url}
\usepackage{hyperref}
\begin{document}
\section{Inline Bibliography}\label{sec:inline}
Classic work~\cite{k_plain} and deep learning~\cite{k_opt_doi} laid foundations,
followed by fast attention~\cite{k_arxiv_pmid} and notes~\cite{k_minimal}.
% \bibitem{fake_comment} G. Ghost. Fake comment entry, 2020.
Inline verbatim \verb|\bibitem{fake_verb}| is ignored.
See Section~\ref{sec:inline}.

\begin{thebibliography}{99}

\bibitem{k_plain}
L. A. Zadeh. Fuzzy sets. \emph{Information and Control}, 8(3), 338--353, 1965.

\bibitem[LeCun et~al.(2015)]{k_opt_doi}
Y. LeCun, Y. Bengio, and G. Hinton, ``Deep learning,'' \emph{Nature}, vol.~521, pp.~436--444, 2015. doi:10.1038/nature14539.

\bibitem[{Dao [et~al.](2022)}]{k_arxiv_pmid}
Tri Dao, Daniel~Y. Fu, Stefano Ermon, Atri Rudra, and Christopher R{\'e}.
\newblock FlashAttention: Fast and memory-efficient exact attention with {IO}-awareness.
\newblock \url{https://arxiv.org/abs/2205.14135}, PMID: 9500320, 2022.

\bibitem{k_minimal}
Working Group Technical Note.

\end{thebibliography}
\end{document}
